% Independent derivation, 2026-09-23. CTD1--30.
\section{The full supported tolerant-circle Hom and its signed divisor receiver}
\label{sec:supported-tolerance-duality}

\paragraph{Sources and scope.}
Alain Connes and Caterina Consani,
\href{https://arxiv.org/abs/2205.01391v2}{\emph{Riemann--Roch for
$\overline{\operatorname{Spec}\mathbb Z}$}}, original author source
\nolinkurl{RR-J-final.tex}, define their signed bounded sections in
Proposition \texttt{setup1}, equation \texttt{kerpsi}; the internal Hom
in Appendix A, Lemma \texttt{catcat}, equation \texttt{catcat1};
the total-distance relations in Proposition \texttt{defnA}, equation
\texttt{inequdef}; and their duality in Proposition \texttt{pontrj}
and Theorem \texttt{serredual}. Appendix B, Proposition
\texttt{proppi0}, defines the actual Dold--Kan-derived kernel and
tolerant cokernel. The complete original source and its version-specific
SHA-256 are retained in the source-use ledger.
We read the original section on cohomology, the complete section
\texttt{sec5}, and Appendices A and B for this calculation.

The later Connes--Consani paper
\href{https://arxiv.org/abs/2306.00456v1}{\emph{Riemann--Roch for the
ring $\mathbb Z$}}, original \nolinkurl{RRZ.tex}, equations
\texttt{xdefn} and \texttt{xdefnr}, uses an all-subset bound.
Its preceding equation \texttt{normdefn} states the distinct signed
bound. We prove the exact comparison of these two bounds below.
The full support carrier is the one of The Clankers,
\href{https://github.com/KokunoYumeto/zeta-function-research-reader/blob/a2851f9eb352e7e4376bc629de86fa4ed9c1c434/workbenches/splitzero-tandem/foundations/split-support-geometry-v11/split_support_geometry_arithmetic_curve_v11.tex#L1226}
{\emph{Split Support Geometry}, version 11,
Definition \texttt{def:lattice-split}}.

The objects calculated here are explicitly specified tolerant
$\Gamma$-objects and their internal Homs. We do not redefine the
authors' $H^1(D)$ by declaring one of these objects to be its
Dold--Kan lift. Neither the divisor $H^1$ nor its lift is identified
with an arithmetic spectral $H^1$.

\subsection{The original divisor and its supported tolerant circle}

Let $L$ be a bounded distributive lattice with $0_L\ne1_L$.
Retain
\begin{equation}\tag{CTD1}
\begin{gathered}
D=\sum_p d_p[p]+d_\infty[\infty],\qquad
N=\prod_p p^{-d_p}>0,\qquad a=e^{d_\infty}>0,\\
I(D)=N\mathbb Z,\qquad
\mathcal O(D)_f=\prod_p p^{-d_p}\mathbb Z_p
=N\widehat{\mathbb Z},\\
A_N=\mathbb R/N\mathbb Z,\qquad
d_N(\bar x,\bar z)=\min_{n\in\mathbb Z}|x-z-Nn|.
\end{gathered}
\end{equation}
The equality $I(D)=N\mathbb Z$ follows by applying the valuation
inequalities to $q/N$. The finite product retains every divisor
coefficient; $N$ need not be an integer.

For an abelian group $A$ define the commutative monoid
\begin{equation}\tag{CTD2}
G_L(A)=\{(0,\lambda):\lambda\in L\}\cup(A\times\{1_L\}),\quad
(x,\lambda)+(z,\nu)=(x+z,\lambda\vee\nu),\quad
\tau=(0,0_L),\quad e=(0,1_L).
\end{equation}
It is a semimodule over the semiring $S=G_L(\mathbb Z)$ by
$(n,\kappa)(x,\lambda)=(nx,\kappa\wedge\lambda)$.
The carrier condition is preserved: a nonzero product amplitude
forces both relevant labels to be top. Distributivity of joins
and meets proves both distributive laws of this scalar action.

For a commutative monoid $M$, write $HM(r_+)=M^r$.
A pointed map acts by adding along each nonbasepoint fibre; the
empty fibre has value $\tau$. This defines the Eilenberg--MacLane
$\Gamma$-object of the monoid. Define the tolerant object
\begin{equation}\tag{CTD3}
\mathcal T^{L,\mathrm{amp}}_{N,a}(r_+)=G_L(A_N)^r,\qquad
(X,Z)\in\mathcal R^{N,a,\mathrm{amp}}_r
\ \Longleftrightarrow\
\sum_{j=1}^r d_N(x_j,z_j)\le a,
\end{equation}
where $x_j,z_j$ are the amplitudes, and all support labels remain
in $X,Z$. The basepoint is the all-$\tau$ tuple.
The relation is reflexive and symmetric. The triangle inequality
gives
\[
\sum_i d_N\!\left(\sum_{f(j)=i}x_j,\sum_{f(j)=i}z_j\right)
\le\sum_jd_N(x_j,z_j),
\]
so every $\Gamma$ structure map preserves the relation. Thus
(CTD3) is an object of the original category $\Gamma\mathcal T_*$.
Its support-sensitive carrier has a tolerance pulled back from
amplitudes; in particular it relates $\tau$ to $e$.
The label-exact object has the same carrier and the smaller relation
\[
(X,Z)\in\mathcal R^{N,a,\mathrm{eq}}_r
\quad\Longleftrightarrow\quad
(X,Z)\in\mathcal R^{N,a,\mathrm{amp}}_r
\quad\hbox{and}\quad \lambda_j=\nu_j\quad(1\le j\le r),
\]
where $\lambda_j,\nu_j$ are the respective labels.
It too is a tolerant $\Gamma$-object: reflexivity and symmetry
hold, and a fold of coordinatewise equal labels produces equal
joins. Denote it by $\mathcal T^{L,\mathrm{eq}}_{N,a}$.
Throughout the principal calculation below,
$\mathcal T^L_{N,a}$ without a superscript means this
\emph{label-exact} object. Both versions and all four directions
of internal-Hom comparison are retained in (CTD27)--(CTD28).

Use the exact original internal-Hom convention
\begin{equation}\tag{CTD4}
\mathfrak H_D(k_+)
=\operatorname{Hom}_{\Gamma\mathcal T_*}
 \bigl(\mathcal T^L_{N,a},
       \mathcal T^L_{1,1/4}(k_+\wedge-)\bigr).
\end{equation}
Its basepoint is the constant all-$\tau$ map. A pointed map
$k_+\to\ell_+$ acts by postcomposition with the corresponding
target fold. We will also calculate the subfunctor
$\mathfrak H_D^S$ whose underlying monoid maps are $S$-linear.
This is an explicitly imposed restriction on morphisms, not a
change to the original category in (CTD4). In particular, fixed
tolerance bounds are not asserted to be preserved by multiplication
by every integer in $S$.

\subsection{All natural maps and the additional support maps}

A natural map $HM\to HN(k_+\wedge-)$ is exactly a $k$-tuple of
pointed additive monoid maps $M\to N$. Indeed, naturality for the
coordinate projections forces the map at level $r_+$ to apply its
level-one maps to each of the $r$ input coordinates. Naturality
for the two-input fold then forces additivity, and pointedness
forces preservation of zero. Conversely these properties commute
with every finite fibre sum, hence give a natural map.

Let
\[
J_0(L)=\{\theta:L\to L:\theta(0_L)=0_L,\
 \theta(\lambda\vee\nu)=\theta(\lambda)\vee\theta(\nu)\},
\qquad
J_{01}(L)=\{\theta\in J_0(L):\theta(1_L)=1_L\}.
\]
These are definitions of sets of maps; $J_0(L)$ here does not
denote the join-irreducible elements of $L$.

Every pointed additive monoid map $F:G_L(A)\to G_L(B)$ has
exactly the form
\begin{equation}\tag{CTD5}
F_{\chi,\theta}(x,\lambda)
  =(\chi(x),\theta(\lambda)),\qquad
\chi\in\operatorname{Hom}_{\mathbb Z}(A,B),\quad
\theta\in J_0(L),\quad
\chi\ne0\Longrightarrow\theta(1_L)=1_L.
\end{equation}
Here is a proof, including the support constraint. The additive
idempotents of $G_L(B)$ are exactly $(0,\mu)$, because $b+b=b$
in the group $B$ implies $b=0$. Thus $F(0,\lambda)=(0,\theta(\lambda))$
with $\theta$ preserving joins and bottom. Set $\mu=\theta(1_L)$.
For $\widehat x=(x,1_L)$, the identities
$\widehat x+e=\widehat x$ and $\widehat x+\widehat{-x}=e$
force the label of $F(\widehat x)$ to be both above and below
$\mu$, hence equal to $\mu$. Write its amplitude as $\chi(x)$.
The group law among the $\widehat x$ proves additivity of $\chi$.
If $\chi(x)\ne0$ for some $x$, the target carrier forces $\mu=1_L$.
This proves necessity and uniqueness. Conversely (CTD5) is
well-defined, preserves bottom, and preserves addition by the
stated properties of $\chi$ and $\theta$.

The $S$-linear maps among (CTD5) are exactly
\begin{equation}\tag{CTD6}
F_{\chi,\mu}(x,\lambda)
=(\chi(x),\lambda\wedge\mu),\qquad
\mu\in L,\qquad \chi\ne0\Longrightarrow\mu=1_L.
\end{equation}
In fact scalar linearity applied to $(0,\lambda)e$ gives
$\theta(\lambda)=\lambda\wedge\theta(1_L)$. Conversely the meet
formula is additive by distributivity of $L$, and
\[
F_{\chi,\mu}((n,\kappa)(x,\lambda))
=(n\chi(x),\kappa\wedge\lambda\wedge\mu)
=(n,\kappa)F_{\chi,\mu}(x,\lambda).
\]
It also preserves $\tau$. Consequently the underlying, unbounded
semimodule Hom is exactly
\[
\operatorname{Hom}_S(G_L(A),G_L(B))
 \cong G_L(\operatorname{Hom}_{\mathbb Z}(A,B)).
\]
A $k$-tuple of such maps has $k$ independent labels $\mu_i$.
It must not be replaced by $G_L(\operatorname{Hom}(A,B)^k)$,
which has a single common label.

At the unbounded level, pointwise addition and scalar action give
\begin{equation}\tag{CTD7}
F_{\chi,\mu}+F_{\psi,\nu}=F_{\chi+\psi,\mu\vee\nu},
\qquad
(n,\kappa)F_{\chi,\mu}=F_{n\chi,\kappa\wedge\mu}.
\end{equation}
In particular a nonzero amplitude map and its negative add to
$F_{0,1_L}$, not to the zero map $F_{0,0_L}$.

\subsection{The all-level tolerance excludes discontinuous characters}

For maps (CTD5) from $G_L(A_N)$ to $G_L(\mathbb R/\mathbb Z)$,
equal input labels give equal output labels for every support
map $\theta_i$. The numerical part of the tolerance is the
amplitude distance. The
necessary and sufficient condition on a $k$-tuple $(\chi_i)$ is
\begin{equation}\tag{CTD8}
\sum_{j=1}^r d_N(x_j,0)\le a
\quad\Longrightarrow\quad
\sum_{j=1}^r\sum_{i=1}^k d_1(\chi_i(x_j),0)\le\frac14
\quad\text{for every }r\ge1.
\end{equation}
Necessity follows by taking the comparison tuple to have zero
amplitudes and exactly the same coordinate labels as the input;
all these zero-amplitude labelled coordinates are admissible.
Conversely use additivity on amplitude differences in any pair;
equal input labels already give equal output labels.
This proves sufficiency at all levels, with every support retained.

No continuity is assumed. Repeat the same $x$ in $m$ input
coordinates. Formula (CTD8) gives
\begin{equation}\tag{CTD9}
d_N(x,0)\le a/m
\quad\Longrightarrow\quad
\sum_{i=1}^k d_1(\chi_i(x),0)\le\frac1{4m}.
\end{equation}
For each positive $\epsilon$, choose $m$ with $1/(4m)<\epsilon$.
Then $d_N(x,0)\le a/m$ forces the distance of each $\chi_i(x)$
from zero to be less than $\epsilon$. Thus each additive map is
continuous at zero and, by translation, uniformly continuous.
This eliminates all discontinuous additive maps by the actual
all-level tolerance, not by an additional hypothesis.

For completeness, every continuous additive
$\chi:A_N\to\mathbb R/\mathbb Z$ is uniquely
\begin{equation}\tag{CTD10}
\chi_t(\bar x)=tx\bmod\mathbb Z,\qquad t\in N^{-1}\mathbb Z.
\end{equation}
To prove this, compose with $\mathbb R\to A_N$ and call the
result $\varphi$. On a sufficiently small interval about zero,
continuity gives a unique continuous real lift $g$ with values
in $(-1/8,1/8)$. Whenever $x,y,x+y$ lie in that interval,
$g(x+y)-g(x)-g(y)$ is an integer of absolute value less than
$3/8$, hence is zero. Extend $g$ to $\mathbb R$ by
$G(x)=2^m g(2^{-m}x)$ for sufficiently large $m$.
Local additivity proves independence of $m$ and global
additivity. The extension agrees with $g$ near zero and is
continuous there; hence $G(x)=tx$, as follows first for rational
$x$ and then for real $x$ by continuity. Addition in the circle
gives $\varphi(x)=tx\bmod\mathbb Z$ at every $x$.
Since $\varphi(N)=0$, $tN\in\mathbb Z$.
Uniqueness follows because $(t-t')x$ cannot be integer for every
sufficiently small real $x$ unless $t=t'$.

The exact coefficient bound is
\begin{equation}\tag{CTD11}
\text{(CTD8)}
\quad\Longleftrightarrow\quad
t_i\in N^{-1}\mathbb Z\ (1\le i\le k),
\qquad a\sum_{i=1}^k|t_i|\le\frac14.
\end{equation}
For sufficiency choose a closest representative $x-Nn$ in
$[-N/2,N/2]$. Since $t_iNn$ is integer,
\[
d_1(t_ix,0)\le |t_i|\,d_N(\bar x,0).
\]
Summing this inequality proves (CTD8).
For necessity choose an integer $m$ so large that
$a/m<N/2$ and $|t_i|a/m<1/2$ for every $i$. Repeat
$\bar x=\overline{a/m}$ exactly $m$ times. Its total source
distance is exactly $a$, and no target term wraps around the
circle. Consequently (CTD8) is precisely
$m\sum_i|t_i|a/m\le1/4$. This proves the displayed bound,
including its boundary equality.

\subsection{The entire internal Hom and its canonical linear retract}

The preceding proofs determine (CTD4) completely:
\begin{equation}\tag{CTD12}
\begin{split}
\mathfrak H_D(k_+)=\{((t_i,\theta_i))_{i=1}^k:\;&
 t_i\in N^{-1}\mathbb Z,\quad
 a\sum_i|t_i|\le1/4,\quad \theta_i\in J_0(L),\\
 &t_i\ne0\Longrightarrow\theta_i\in J_{01}(L)\}.
\end{split}
\end{equation}
The actual natural transformation represented by this tuple acts
at level $r_+$ by
\[
((x_j,\lambda_j))_{j=1}^r
\longmapsto
((t_i x_j\bmod\mathbb Z,\theta_i(\lambda_j)))_
 {1\le i\le k,\ 1\le j\le r}.
\]
A pointed map $\alpha:k_+\to\ell_+$ acts by
\begin{equation}\tag{CTD13}
(t_i,\theta_i)_i\longmapsto
\left(\sum_{\alpha(i)=h}t_i,\
      \bigvee_{\alpha(i)=h}\theta_i\right)_{h=1}^{\ell},
\end{equation}
where the join of maps is pointwise and an empty join is the
constant-bottom map. It preserves (CTD12), since the triangle
inequality preserves the coefficient bound, a pointwise join of
join-preserving maps preserves joins, and a nonzero resulting
coefficient has some nonzero input coefficient and therefore
top output at $1_L$. Formula (CTD13) also follows directly by
folding the displayed natural transformation.

Its linear subfunctor is exactly
\begin{equation}\tag{CTD14}
\mathfrak H_D^S(k_+)=
\{((t_i,\mu_i))_{i=1}^k:
t_i\in N^{-1}\mathbb Z,\ a\sum_i|t_i|\le1/4,\
\mu_i\in L,\ t_i\ne0\Rightarrow\mu_i=1_L\}.
\end{equation}
The maps
\begin{equation}\tag{CTD15}
\begin{gathered}
\iota_D:\mathfrak H_D^S\longrightarrow\mathfrak H_D,\qquad
(t_i,\mu_i)\longmapsto(t_i,[\lambda\mapsto\lambda\wedge\mu_i]),\\
\rho_D:\mathfrak H_D\longrightarrow\mathfrak H_D^S,\qquad
(t_i,\theta_i)\longmapsto(t_i,\theta_i(1_L)),\qquad
\rho_D\iota_D=\mathrm{id}
\end{gathered}
\end{equation}
are pointed natural transformations. For the section, naturality
uses $\lambda\wedge(\bigvee_i\mu_i)
=\bigvee_i(\lambda\wedge\mu_i)$; for the retraction it uses
$(\bigvee_i\theta_i)(1_L)=\bigvee_i\theta_i(1_L)$.
The fibre of $\rho_D$ over a tuple $(t_i,\mu_i)$ is exactly
\[
\prod_{i=1}^k\{\theta\in J_0(L):\theta(1_L)=\mu_i\}.
\]
All these fibres are retained in (CTD12). The pointed kernel
of $\rho_D$ is a singleton: if $\theta(1_L)=0_L$, monotonicity
forces $\theta$ to be the constant-bottom map. This does not
make $\rho_D$ injective.

For example take $L=\{0_L<u<1_L\}$. The map
$\theta(0_L)=0_L$, $\theta(u)=0_L$, $\theta(1_L)=1_L$
preserves joins and bottom, but differs from the identity
meet map with $\mu=1_L$. Already the zero character with this
$\theta$ and the zero character with the identity map are
distinct elements of $\mathfrak H_D(1_+)$ having the same
image under $\rho_D$, for every $D$.
Thus the proposed full supported coefficient carrier is a
canonical retract, not the entire original-category Hom in
general. For the two-element lattice all members of $J_0(L)$
are either the bottom map or the identity, so $\rho_D$ is an
isomorphism in that particular case.

\subsection{The canonical divisor and the exact signed section map}

Keep the source's actual divisor $K=-2[2]$. Then
\begin{equation}\tag{CTD16}
\begin{gathered}
I(K)=4\mathbb Z,\qquad e^{K_\infty}=1,\qquad
\mathcal T_K^L=\mathcal T^L_{4,1},\\
\mathcal T^L_{4,1}\xrightarrow{\ \cong\ }\mathcal T^L_{1,1/4},
\qquad (z\bmod4\mathbb Z,\lambda)
\longmapsto(z/4\bmod\mathbb Z,\lambda),\\
I(K-D)=\frac4N\mathbb Z,\qquad
e^{(K-D)_\infty}=a^{-1}.
\end{gathered}
\end{equation}
The circle map and its inverse multiplication by $4$ preserve
every label. The distance identity $d_1(z/4,w/4)=d_4(z,w)/4$
proves preservation of the displayed tolerances at every level.

Define the supported signed section module by its full levels:
\begin{equation}\tag{CTD17}
\mathcal H^{0,\pm}_L(K-D)(k_+)
=\{((q_i,\mu_i))_{i=1}^k:
q_i\in(4/N)\mathbb Z,\ q_i\ne0\Rightarrow\mu_i=1_L,\
\sum_i|q_i|\le1/a\}.
\end{equation}
Its folds sum amplitudes and join labels. The triangle inequality
proves that each fold preserves the bound. This is the
support lift of the signed section module of the 2022 source,
with its full original ideal and real bound.

The exact pointed natural isomorphism is
\begin{equation}\tag{CTD18}
\mathfrak H_D^S\xrightarrow{\ \cong\ }\mathcal H^{0,\pm}_L(K-D),
\qquad ((t_i,\mu_i))\longmapsto((4t_i,\mu_i)),
\qquad q_i\longmapsto t_i=q_i/4.
\end{equation}
The ideal and the bound follow respectively from
$t_i\in N^{-1}\mathbb Z$ and
$a\sum_i|t_i|\le1/4$. All folds commute with multiplication by
$4$, and all labels are unchanged. Combining (CTD15) and
(CTD18) therefore gives an explicit retract of the complete
internal Hom onto the signed supported section module.

The underlying evaluation before the circle comparison is
\begin{equation}\tag{CTD19}
((q,\mu),(x\bmod N\mathbb Z,\lambda))
\longmapsto(qx\bmod4\mathbb Z,\mu\wedge\lambda).
\end{equation}
It is well-defined because $qN\in4\mathbb Z$.
For a $k$-tuple of coefficients and an $r$-tuple of inputs,
its $(i,j)$ entry is the same expression with $q_i,x_j$.
The bound $\sum_i|q_i|\le1/a$ implies
\[
\sum_{i,j}d_4(q_ix_j,q_iz_j)
\le\left(\sum_i|q_i|\right)\sum_jd_N(x_j,z_j)\le1.
\]
Thus evaluation gives the actual required tolerant morphism into
$\mathcal T_K^L(k_+\wedge-)$, with the original factor $4$
and all support meets present.

For principal covariance retain
$\operatorname{div}(c)=\sum_pv_p(c)[p]-\log|c|[\infty]$.
For $D'=D+\operatorname{div}(c)$ one has
$N'=N/|c|$ and $a'=a/|c|$. The exact isomorphism
$\mathcal T_D^L\to\mathcal T_{D'}^L$ is
$(x,\lambda)\mapsto(x/c,\lambda)$.
Its inverse induces on internal Homs, and on their signed
receivers, the maps
\begin{equation}\tag{CTD20}
(t_i,\theta_i)\longmapsto(ct_i,\theta_i),\qquad
(q_i,\mu_i)\longmapsto(cq_i,\mu_i).
\end{equation}
Indeed $a'\sum|ct_i|=a\sum|t_i|$, and
$I(K-D')=(4|c|/N)\mathbb Z=c(4/N)\mathbb Z$.
All inverse formulas divide by $c$; all support labels remain.
Negative $c$ reverses the amplitude signs and leaves the absolute
bounds and support joins unchanged.

\subsection{The strongest direct comparison with the all-subset module}

Write $c_D=1/a$ and $M_D=4/N$, retaining these exact values.
Let $\mathcal U_L(M_D,c_D)$ be the RDT all-subset module with
amplitudes in $M_D\mathbb Z$ and every subset sum in
$[-c_D,c_D]$. Let $\mathcal S_L(M_D,c_D)$ denote the signed
module (CTD17). For a tuple $q$, put
\[
P(q)=\sum_{q_i>0}q_i,\qquad
Q(q)=\sum_{q_i<0}|q_i|.
\]
The original conditions, with every sign retained, are
\begin{equation}\tag{CTD21}
\begin{array}{ll}
q\in\mathcal S_L(M_D,c_D)&\Longleftrightarrow\
P(q)+Q(q)\le c_D,\\
q\in\mathcal U_L(M_D,c_D)&\Longleftrightarrow\
P(q)\le c_D\ \hbox{and}\ Q(q)\le c_D.
\end{array}
\end{equation}
All support conditions in both lines are identical. The identity
on amplitudes and labels gives the exact natural inclusions
\begin{equation}\tag{CTD22}
\mathcal H^{0,\pm}_L(K-D)
=\mathcal S_L(M_D,c_D)
\ \lhook\joinrel\longrightarrow\
\mathcal H^0_L(K-D)=\mathcal U_L(M_D,c_D)
\ \lhook\joinrel\longrightarrow\
\mathcal S_L(M_D,2c_D).
\end{equation}
The second inclusion follows from $P+Q\le2c_D$.
These are actual componentwise maps, commuting with all folds.

They agree at level $1_+$, but the first inclusion is strict
at level $2_+$ whenever $m=\lfloor c_D/M_D\rfloor\ge1$:
take $q=mM_D$ and the full supported tuple
$((q,1_L),(-q,1_L))$. It is admitted by the all-subset condition,
but $2mM_D>c_D$ because $c_D<(m+1)M_D\le2mM_D$.
It therefore fails the signed condition. If $m=0$, both modules
are precisely $H(L,\vee,0_L)$ in all levels.

There is no natural transformation from the middle module in
(CTD22) to the first whose level-one map is the identity when
$m\ge1$. Projection naturality would force its value at the
displayed two-coordinate tuple to have those same coordinates,
which are not an element of the required signed target.
This proves the exact limitation; it does not erase the
inclusions or the full retained objects.

The exact excess object is the pointed $\Gamma$-cofiber
\begin{equation}\tag{CTD23}
\mathcal Q_{D,L}(k_+)
=\mathcal U_L(M_D,c_D)(k_+)/
       \mathcal S_L(M_D,c_D)(k_+).
\end{equation}
Here the indicated subset is collapsed to one basepoint.
The remaining elements are precisely the full labelled tuples
with $P,Q\le c_D$ and $P+Q>c_D$.
Folds are the original folds, followed by this stated collapse;
well-definedness follows because the signed object is a
subfunctor. This is a pointed-set cofiber, not a claimed
abelian quotient or a faithful replacement for the two original
modules. Equations (CTD12), (CTD15), (CTD18), and (CTD22)
retain every source and comparison in the duality calculation.

\subsection{Pullback along the actual adelic quotient}

The source's Proposition \texttt{thedimh1} uses the adelic
quotient with finite subgroup $\mathcal O(D)_f$, not just an
unqualified real coordinate projection. Explicitly, for
$x=(x_\infty,x_f)\in\mathbb A$, choose $r\in\mathbb Q$ such that
$x_f-r\in\mathcal O(D)_f$ and set
\begin{equation}\tag{CTD24}
\pi_D(x)=x_\infty-r\bmod N\mathbb Z.
\end{equation}
Such an $r$ exists by the following complete finite calculation.
Choose a positive integer $M$ with $M x_f/N\in\widehat{\mathbb Z}$.
For each prime dividing $M$, prescribe the residue of an integer
$j$ to be that of $M x_p/N$ modulo $p^{v_p(M)}$.
The Chinese remainder theorem supplies $j$. At primes not
dividing $M$, both $M x_p/N$ and $j$ are integral and $M$ is a
unit. Hence $M x_f/N-j\in M\widehat{\mathbb Z}$ at every prime.
Take $r=Nj/M\in\mathbb Q$; then
$x_f-r\in N\widehat{\mathbb Z}$, as required. Two choices differ by
$\mathbb Q\cap N\widehat{\mathbb Z}=N\mathbb Z$, so the
formula is independent of the choice. Choosing sums of the
rational representatives proves additivity. It is surjective
by restricting to adeles $(t,0)$, and
\[
\ker\pi_D=\mathbb Q+
 \bigl(\{0\}\times\mathcal O(D)_f\bigr).
\]
Both containments follow directly from (CTD24).

Give $H(G_L(\mathbb A))$ the relation
$\sum_jd_N(\pi_Dx_j,\pi_Dz_j)\le a$ together with
equality of the two labels in each coordinate,
and denote the resulting explicit tolerant object by
$\widetilde{\mathcal T}_D^L$. Then precomposition with the
label-preserving $G_L(\pi_D)$ gives exact natural isomorphisms
\begin{equation}\tag{CTD25}
\begin{gathered}
\operatorname{Hom}_{\Gamma\mathcal T_*}
(\widetilde{\mathcal T}_D^L,\mathcal T^L_{1,1/4}(k_+\wedge-))
\cong\mathfrak H_D(k_+),\\
\operatorname{Hom}^{S}_{\Gamma\mathcal T_*}
(\widetilde{\mathcal T}_D^L,\mathcal T^L_{1,1/4}(k_+\wedge-))
\cong\mathfrak H_D^S(k_+).
\end{gathered}
\end{equation}
For the proof apply (CTD5) to each component.
If $x\in\ker\pi_D$, compare repeated inputs $(x,1_L)$ with
the same number of coordinates $(0,1_L)$. They have equal
labels and zero source distance. Preservation of the target tolerance
forces $m\,d_1(\chi_i(x),0)\le1/4$ for every $m$, hence
$\chi_i(x)=0$. Thus each amplitude homomorphism factors uniquely
through $\pi_D$. All support maps $\theta_i$ are unchanged.
Surjectivity of $\pi_D$ shows that the factored maps satisfy
exactly the circle tolerance. Conversely composition with
$\pi_D$ preserves it by definition. This proves both
isomorphisms and compatibility with every internal-Hom fold.

This pullback has an exact receiving map from the native supported
adelic construction: (ASC7) and (ASC17)--(ASC18) identify its relation
with the maximal symmetric subrelation of the degree-zero directed
relation of $X(D)$. Thus (CTD25) applies to that calculated symmetric
part, with its original vertex carrier and all its fibres. It is not
an identification of the directed relation with a symmetric relation,
nor a claim that (ASC8)'s pointed spherical set is the signed section
module. The actual signed sections are the retract of the complete
relative zero-section object in (ASC9)--(ASC13).

The assertion is a proof about the explicitly defined pullback
object. In particular the functor $G_L$ does not preserve
pointed kernels of group maps. For an additive map $f:A\to B$,
\begin{equation}\tag{CTD26}
G_L(f)(x,\lambda)=(f(x),\lambda),\qquad
G_L(f)^{-1}(\tau_B)=\{\tau_A\},\qquad
(\mathrm{amplitude}\circ G_L(f))^{-1}(0)=G_L(\ker f).
\end{equation}
The first kernel statement follows because an output equal to
$\tau_B$ requires $\lambda=0_L$, hence $x=0$ by the carrier
condition. The amplitude-zero fibre follows by requiring only
$f(x)=0$. This exact difference prevents an unproved transfer
of the original Dold--Kan kernel through $G_L$.
The original source's $H^1(D)$ and its amplitude comparison are
retained; the support lift calculated here has been specified
and its actual maps proved, without identifying it with an
unconstructed spectral cohomology.

\subsection{Both tolerances and all four directions of comparison}

The identity on the full carrier induces a natural morphism
\begin{equation}\tag{CTD27}
j_{N,a}:\mathcal T^{L,\mathrm{eq}}_{N,a}
\longrightarrow\mathcal T^{L,\mathrm{amp}}_{N,a}.
\end{equation}
It retains every element and label, and preserves relations
because the label-exact relation is contained in the
amplitude relation. The reverse identity is not a tolerant
morphism for a nontrivial lattice: $\tau$ and $e$ are related
in the amplitude object and are not related in the label-exact
one. The map (CTD27) forgets a condition on pairs; it does
not identify those two points.

For superscripts $\alpha,\beta\in\{\mathrm{eq},\mathrm{amp}\}$,
write $\mathfrak H_D^{\alpha\to\beta}$ for the internal Hom
from $\mathcal T^{L,\alpha}_{N,a}$ to
$\mathcal T^{L,\beta}_{1,1/4}$. Then the exact classifications,
as pointed $\Gamma$-objects, are
\begin{equation}\tag{CTD28}
\begin{gathered}
\mathfrak H_D^{\mathrm{eq}\to\mathrm{eq}}
\ \cong\
\mathfrak H_D^{\mathrm{amp}\to\mathrm{amp}}
\ \cong\
\mathfrak H_D^{\mathrm{eq}\to\mathrm{amp}}
\ =\ \hbox{the object in (CTD12)},\\
\mathfrak H_D^{\mathrm{amp}\to\mathrm{eq}}=\{*\}.
\end{gathered}
\end{equation}
The same statement holds for the $S$-linear subfunctors, with
(CTD14) as the common nontrivial object.
Here the displayed isomorphisms keep the actual natural
transformations, their characters, and all their support maps.
They are not isomorphisms between the original tolerant
source objects themselves.

To prove the first three classifications, a morphism in any
of those directions is still of the form (CTD5) by
$\Gamma$-naturality. In each case the repeated top-labelled
inputs and top-labelled zero comparisons used in
(CTD8)--(CTD11) are related in the source, so they force the
same continuity and exact coefficient bound. Conversely
that coefficient bound preserves the amplitude part of the
relation. If the target relation is label-exact in one of
these three cases, the source relation is also label-exact,
and applying $\theta_i$ to equal labels gives equal labels.
This proves necessity and sufficiency. Precomposing or
postcomposing with (CTD27) realizes the displayed
isomorphisms on actual maps.

For the remaining direction, every pair
$((0,\lambda),\tau)$ is related in the source amplitude
tolerance. A morphism into the label-exact target must
therefore satisfy $\theta_i(\lambda)=\theta_i(0_L)=0_L$
for every $\lambda$ and every output coordinate $i$.
Thus every $\theta_i$ is the bottom map. Since
$\theta_i(1_L)=0_L\ne1_L$, (CTD5) forces $\chi_i=0$.
The only possible natural transformation is the constant
all-$\tau$ transformation, which is indeed a morphism.
This proves the final assertion at every level.

Replacing the label-exact adelic relation in (CTD25) by its
amplitude version gives the corresponding same-tolerance
pullback isomorphism, by the identical repeated-input proof.
The mixed direction with amplitude source and label-exact
target remains the one-point object. Consequently there is
no silent exchange of support tolerances in the duality
diagram: the label-exact calculation has its complete
nontrivial character and support receiver, the amplitude
calculation has the same internal-Hom parameters when used
on both sides, and switching only the target to label-exact
destroys all nonzero maps.

\subsection{Every finite-support cardinality and the zero-character range}

For finite $L$, set
$j_0=\#J_0(L)$, $j_{01}=\#J_{01}(L)$, and
$m=\lfloor N/(4a)\rfloor$. Then the complete level counts are
\begin{equation}\tag{CTD29}
\begin{gathered}
\#\mathfrak H_D(k_+)
=\sum_{r=0}^{k}\binom{k}{r}2^r\binom{m}{r}
       j_{01}^{\,r}j_0^{\,k-r},\\
\#\mathfrak H_D^S(k_+)
=\sum_{r=0}^{k}\binom{k}{r}2^r\binom{m}{r}|L|^{\,k-r}.
\end{gathered}
\end{equation}
Use $\binom m0=1$ and $\binom mr=0$ for $r>m$.
Indeed write each character coefficient as $t_i=n_i/N$,
retaining the original $N$. The exact bound (CTD11) is
$\sum_i|n_i|\le N/(4a)$, hence $\sum_i|n_i|\le m$.
Choose its $r$ nonzero positions, their $2^r$ signs, and their
ordered positive magnitudes with total at most $m$.
The latter have number $\binom mr$: their successive partial
sums form an increasing $r$-element subset of
$\{1,\ldots,m\}$, with inverse successive differences.
Each nonzero coefficient allows exactly $j_{01}$ support
maps in the full Hom, and each zero coefficient allows
exactly $j_0$. In the linear subfunctor the respective numbers
are one and $|L|$. This proves both formulas, including
the single empty tuple when $k=0$.

Without any finiteness hypothesis on $L$, if $a>N/4$ then
$m=0$ and the exact entire $\Gamma$-objects are
\begin{equation}\tag{CTD30}
\mathfrak H_D=H(J_0(L),\vee,0),\qquad
\mathfrak H_D^S=H(L,\vee,0_L)\qquad(a>N/4).
\end{equation}
The join of maps in the first expression is pointwise and
its zero is the constant-bottom map. All character
coefficients vanish by (CTD11), but their support maps do
not disappear. For example the map with zero character
and $\theta=\mathrm{id}_L$ sends every top-supported input
to $e$ and every $(0,\lambda)$ to itself. It differs from
the constant-$\tau$ map and satisfies the stated tolerances.
Equation (CTD30) follows also directly from the fold formula
(CTD13); the linear retract is exactly evaluation at $1_L$
and its meet-map section. Thus absence of nonzero amplitude
characters leaves the entire calculated support receiver.

